\documentclass[11pt,a4paper]{article}
\usepackage[margin=2cm,headheight=14pt]{geometry}
\usepackage[T1]{fontenc}
\usepackage[utf8]{inputenc}
\usepackage{lmodern,microtype,amsmath,booktabs,tabularx,array,enumitem,fancyhdr}
\usepackage{xcolor,listings,textcomp}
\usepackage[colorlinks=true,urlcolor=blue,linkcolor=blue]{hyperref}
\hypersetup{pdftitle={Research Methods 2026 -- Coding Assignment 1},pdfauthor={Swapnil Singh}}
\pagestyle{fancy}
\fancyhf{}
\fancyhead[L]{Research Methods 2026}
\fancyhead[R]{Coding Assignment 1}
\fancyfoot[C]{\thepage}
\setlength{\parindent}{0pt}
\setlength{\parskip}{5pt}
\setlist{topsep=3pt,itemsep=2pt,parsep=0pt,leftmargin=*}
\renewcommand{\arraystretch}{1.12}
\newcommand{\py}[1]{\texttt{\detokenize{#1}}}
\newcommand{\question}[2]{\section*{#1\hfill{\normalsize[#2 marks]}}}
\lstset{language=Python,basicstyle=\ttfamily\small,keywordstyle=\color{blue!65!black},
  commentstyle=\color{black!60},stringstyle=\color{red!55!black},
  backgroundcolor=\color{black!3},frame=single,rulecolor=\color{black!15},
  numbers=left,numberstyle=\tiny\color{black!50},numbersep=8pt,
  xleftmargin=1.8em,framexleftmargin=0.5em,breaklines=true,
  showstringspaces=false,upquote=true,columns=fullflexible,keepspaces=true,
  aboveskip=6pt,belowskip=6pt}

\begin{document}
\begin{center}
{\Large\textbf{Coding Assignment 1}}\\[4pt]
{\large Household support: implement, inspect, and verify}\\[5pt]
Research Methods 2026 \quad | \quad Swapnil Singh
\end{center}

\begin{tabularx}{\linewidth}{@{}lX@{}}
\textbf{Deadline} & Monday 9 November 2026, 23:59 (Europe/Vilnius) \\
\textbf{Weight} & 100 assignment marks; contributes 10\% of the whole course grade \\
\textbf{Work} & Individual; plan approximately three hours \\
\textbf{Hand in} & \py{exercise_01.ipynb} and \py{README.md} through a GitHub pull request \\
\end{tabularx}

\textbf{Scope.} Use the material from Lectures 1--4: running a project, basic
Python, conditions, loops, functions, inline functions (\py{lambda}), and checks.
All inputs are Python lists supplied below. No CSV import, pandas, NumPy,
plotting, or external data is required. Git commands and version-control
knowledge are not assessed.
Use built-in Python tools
and, for floating-point checks, \py{math.isclose}. List comprehensions and
\py{try}/\py{except} are not required.

\textbf{AI assistance is permitted.}

\section*{Inputs and policy specification}
These eight households are invented for this assignment. Income and consumption
are EUR per month; household size is persons. Positions in the four lists refer
to the same household. Preserve the supplied lists.

\begin{lstlisting}
household_ids = ["A01", "A02", "A03", "A04",
                 "A05", "A06", "A07", "A08"]
income_text = ["1200", "0", "1500", "700",
               "2400", "800", "0", "1800"]
consumptions = [1250, 150, 1600, 850, 2600, 950, 0, 1400]
household_sizes = [2, 1, 3, 1, 4, 2, 2, 3]
\end{lstlisting}

For income $I$, consumption $C$, and size $n$, the hypothetical monthly transfer is
\[
T(I,C,n)=
\begin{cases}
0, & I/n>600,\\
\min\{100n,\;\max\{C-I,0\}\}, & I/n\leq600.
\end{cases}
\]
Thus eligibility includes exactly EUR 600 per person. An eligible household
receives its consumption shortfall, capped at EUR 100 per person. A household
without a shortfall receives zero. Consumption is held fixed throughout.

Define income after support as $Y=I+T$, saving before as $S=I-C$, and saving
after as $S'=Y-C$. A saving share is saving divided by the corresponding income;
it is \py{None} when that income is zero. Negative saving is valid.

\newpage
\question{1. Trace values and notebook state}{12}
The following snippet is separate from the eight-household inputs. Before
running it, write your predictions in a Markdown cell. Keep that first attempt,
then run the snippet and explain any corrections. Initial prediction errors do
not themselves lose marks if the final reasoning is correct.

\begin{lstlisting}
incomes = [1000, 800, 0]
saved_incomes = incomes
first_share = (incomes[0] - 800) / incomes[0]
incomes[0] = 1200
saved_incomes.append(600)
print(incomes)
print(saved_incomes[-1], len(incomes))
print(f"{first_share:.0%}")
\end{lstlisting}

\begin{enumerate}[label=(\alph*)]
\item Give the three output lines and trace the relevant values. \hfill [4]
\item Explain why appending through \py{saved_incomes} affects \py{incomes},
and why changing an income does not update \py{first_share}. \hfill [4]
\item Make \py{saved_incomes} an independent copy of the original list and
recalculate the share after changing income to 1200. Show your changes and the
new three output lines. \hfill [4]
\end{enumerate}

\question{2. Implement the policy as functions}{28}
Write the following three functions. Return values rather than only printing
them. Assume amounts are finite real numbers and household sizes are integers;
you do not need general type, NaN, or infinity checks.

\begin{enumerate}[label=(\alph*)]
\item \py{saving_share(income, consumption)} \hfill [8]

Raise \py{ValueError} if either amount is negative. Otherwise return \py{None}
for zero income, and $(I-C)/I$ for positive income. Check invalid amounts
\emph{before} the zero-income case. Return a fractional share, not a percentage.

\item \py{transfer(income, consumption, size)} \hfill [12]

Raise \py{ValueError} for negative income, negative consumption, or
\py{size <= 0}, before calculating eligibility. For valid inputs, implement
the policy on page 1 exactly. Every valid call must return a number, including
zero when no transfer is payable.

\item \py{post_policy_share(income, consumption, size)} \hfill [8]

Call your \py{transfer} function, add its result to original income, and call
\py{saving_share} using that new income and unchanged consumption. Return the
result. Reuse both helpers; do not duplicate their policy or share formulas.
\end{enumerate}

\textbf{Explain your design in at most 100 words:} distinguish parameters from
arguments, explain why \py{return} is needed, and explain why input checks
precede division. Functions must use their arguments rather than depend on
remembered notebook values. Place these explanations beside the relevant code.
The marks above cover both implementation and explanation.

\newpage
\question{3. Build and interpret a checked report}{30}
Use your functions from Question 2. Write explicit loops where requested so that
your reasoning is visible. Keep full precision in calculations and round only
for display. Counts must identify their denominator or population.

\begin{enumerate}[label=(\alph*)]
\item \textbf{Prepare inputs.} Convert \py{income_text} to a new list of numbers
using a loop; leave the original strings unchanged. Check that all four input
lists have equal lengths \emph{before} processing. Raise \py{ValueError} for a
length mismatch; do not let \py{zip} silently omit households. \hfill [4]

\item \textbf{Calculate household results.} Use an explicit loop over the
aligned inputs. Build result lists and print a labelled report with one line
per household containing: ID, original income, consumption, size, transfer,
saving before, saving after, share before, and share after. Use two decimal
places for EUR amounts and percentage displays; print \texttt{undefined} for
\py{None}. Do not replace an undefined share with zero. \hfill [8]

\item \textbf{Compare totals and averages.} Calculate and report: \hfill [8]
\begin{itemize}
\item total transfer, number of households receiving a positive transfer, and
numbers with negative saving before and after;
\item the pooled saving share before, $\sum_i S_i/\sum_i I_i$, and after,
$\sum_i S'_i/\sum_i Y_i$;
\item the arithmetic mean of defined household saving shares before and after,
with the number included in each mean;
\item the mean after support restricted to households whose share was already
defined \emph{before} support, using the same households for both periods.
\end{itemize}
In at most 120 words, compare the pooled shares and household means. Identify
any change in which households enter the means and explain why comparing means
over different sets can be misleading. No statistical method beyond totals,
counts, and arithmetic averages is needed.

\item \textbf{Use an inline function.} Build a new list containing only the
defined \emph{pre-policy} household shares. Use \py{sorted} with a \py{lambda}
key to order them by distance from a target saving share of \py{0.10}, nearest
first. Print this list and explain what the lambda receives and returns.
Retain the original share values, including their signs; do not output the
distances instead. Preserve the original result list, including \py{None}.
\hfill [4]

\item \textbf{Test a counterfactual.} Make an independent copy of the numeric
income list and change only A01's income from 1200 to 1202. Recompute A01's
transfer, income after support, and saving after support. Show these quantities
under both inputs and the change in each. Explain in at most 80 words whether
higher original income necessarily gives higher income after support under
this policy. Verify that your original input list remains unchanged. \hfill [6]
\end{enumerate}

\textbf{Interpretation matters.} Distinguish a change in a household's saving
amount from a change in its saving share. A program that runs or prints a
plausible number does not establish that it answers the stated question.

\newpage
\question{4. Audit an agent-written function}{20}
An agent proposes the function below. Its intended contract is to return the
arithmetic mean of \emph{defined} saving shares, including valid zero and
negative shares. It must raise \py{ValueError} for unequal list lengths and
return \py{None} when no share is defined, including empty inputs. It must not
modify its inputs. Assume each individual input pair meets Question 2's
numeric and non-negative conditions. \py{saving_share} is your correct helper.

\begin{lstlisting}
def proposed_mean(incomes, consumptions):
    total = 0
    count = 0
    for income, consumption in zip(incomes, consumptions):
        share = saving_share(income, consumption)
        if share is not None and share != 0:
            total = 0
            total = total + share
            count = count + 1
    return total / count
\end{lstlisting}

\begin{enumerate}[label=(\alph*)]
\item Identify \textbf{four distinct defects} and connect each to the intended
contract. Do not count the same defect twice. \hfill [8]
\item Write a repaired \py{mean_defined_share(incomes, consumptions)} using a
loop and your helper. Preserve the intended contract. \hfill [6]
\item Give four small diagnostic input pairs. For each, state the expected
correct outcome, what the proposal actually does, and which defect it exposes.
Use cases that separate the defects as far as possible. An outcome may be a
number, \py{None}, or an exception type. \hfill [6]
\end{enumerate}

Put deliberately failing calls and their observed errors in Markdown code
blocks in the final notebook. You may investigate them in scratch cells first;
you do not need to learn \py{try}/\py{except} to finish this task.

\question{5. Supply evidence and a reproducible hand-in}{10}
\begin{enumerate}[label=(\alph*)]
\item Include at least \textbf{eight meaningful, passing assertions} covering
ordinary saving, negative saving, zero income, the eligibility boundary,
the transfer cap, no shortfall, all-undefined shares, and the counterfactual's
preserved original input. Use \py{is None} for undefined results and
\py{math.isclose} for floating-point comparisons. For four assertions, explain
in one sentence which plausible mistake each would detect. \hfill [6]
\item Restart the kernel and run all cells in order. The final notebook must
run without manual intervention or unhandled exceptions, and reproduce all
reported results. Keep the original first prediction from Question 1 in
Markdown. \hfill [2]
\item In \py{README.md}, give the Python version, instructions to open and run
the notebook, and a brief AI-use record (or a no-AI statement). Submit the
notebook and README through a pull request. \hfill [2]
\end{enumerate}

\end{document}
